\section{Past Actions}

\emph{\textbf{Article XXI of the Ggngral Aereement on Tariffs and Tradg (GATT) 1994 (Security Exceptions), Article XIV bis of the General Aereement on Trade in Sgrvicgs (GATS), and Article Y3 of the WTO Aereement on Trade-Related Aspects of Intellectual Property Riehts (TRIPS):}} Together, these three instruments form a comprehensive \emph{security exception triad} within the framework of the WTO. Article XXI of the GATT is the core legal mechanism within the WTO framework that addresses the friction in balancing free trade and national security. Its provisions allow member states to bypass its general trade obligations if it deems such a step to be necessary to protect its “essential security interests”.\footnote{WTO. (1994). Security Exceptions (Article XXI). General Agreement on Tariffs and Trade (GATT) 1994. World Trade Organization. \url{https://www.wto.org/english/docs_e/legal_e/gatt47_e.htm} art21} Similarly, GATS Article XIV bis and TRIPS Agreement Article 73 are parallel instruments which extend the national security exception to the fields of international services and intellectual property. The former lets members restrict foreign service providers from operating within their borders if considered necessary for essential security, while the latter permits states to sidestep intellectual property protections during situations alike, which is progressively relevant in modern geopolitical tensions where vital dual-use technologies, telecommunications infrastructure, and sophisticated semiconductors are extensively patent-reliant. \footnote{WTO. (1994a). Security Exceptions (Article XIV bis). General Agreement on Trades in Services. World Trade Organization. \url{https://www.wto.org/english/docs_e/legal_e/gats_e.htm} art14\_bis} \footnote{WTO. (1994b). Security Exceptions (Article 73). Agreement on Trade-Related Aspects of Intellectual Property Rights.}

Fundamentally, by establishing legal boundaries that substantially obstruct the total weaponization of global commerce, this triad balances these security imperatives with the continuance of the multilateral trading system. Recent WTO dispute settlement jurisprudence (e.g. \emph{DS512: Russia - Measures Concerning Traffic in Transit}, \emph{DS567: Saudi Arabia - Measures Concerning the Protection of Intellectual Property Rights}, and \emph{DS564: United States - Certain Measures on Steel and Aluminium Products}) has elucidated that while these provisions offer wide discretion, they are not to be exploited as means for economic protectionism. That is, the invocation of a security exception remains subject to objective review to substantiate the existence of a genuine “emergency in international relations”. At times of geopolitical unrest, such a standard of accountability acts as the paramount balancing mechanism. It assists legitimate sovereign defence measures besides preventing major powers from indiscriminately covering routine trade wars and economic nationalism under the guise of national security, consequently preserving the continuity and predictability of the broader global free market.

\emph{\textbf{GATT Article V (Frggdom of Transit) and the WTO Aereement on Trade Facilitation (TFA):}} The freedom of transit for goods moving through the domain of a WTO member is ensured by the GATT Article V. When geopolitical tensions persist, states may impose jarring sanctions on rival nations. However, Article V guarantees that a transit country cannot legally blockade or levy severe tolls on the commerce of third-party nations solely because they

World Trade Organization. \url{https://www.wto.org/english/docs_e/legal_e/31bis_trips_09_e.htm}

\href{https://www.wto.org/english/docs_e/legal_e/31bis_trips_09_e.htm}{\_e/31bis\_trips\_09\_e.htm} are passing through its borders.\footnote{WTO. (1994c). Freedom of Transit (Article V). General Agreement on Tariffs and Trade (GATT) 1994. World Trade Organization. \url{https://www.wto.org/english/docs_e/legal_e/gatt47_e.htm} art5} Furthermore, The WTO’s Agreement on Trade Facilitation (TFA) reinforces this by mandating accelerated customs procedures and transparency, indifferent to the prevalent political climate.\footnote{WTO. (2014). Agreement on Trade Facilitation (Document WT/L/940). World Trade Organization.}

These mechanisms collectively stop geopolitical discord from resulting in precarious logistical paralysis, and help sustain the equilibrium between a country’s sovereign border security and the international community’s access to open, predictable supply chain corridors.

\emph{\textbf{Article 12 of the WTO Aereement on Aericulture:}} States often resort to abrupt export bans to secure their domestic food supply during global crises, which acts as a catalyst to intense geopolitical tension and collapse of international agricultural trade. To confine this protectionist impulse, Article 12 of the WTO Agreement on Agriculture establishes thorough disciplines on export prohibitions and restrictions for foodstuffs.

The mechanism necessitates that when enforcing an export restriction to address a domestic critical shortage, a member state gives due consideration to the effects of such restriction on the importing  members’  food  security.\footnote{WTO. (1994d). Disciplines on Export Prohibitions and Restrictions (Article 12). Agreement on Agriculture. World Trade Organization. \url{https://www.wto.org/english/docs_e/legal_e/ag_e.htm} art12}

\url{https://docs.wto.org/dol2fe/Pages/SS/directdoc.aspx?filename=q\%3A/WT/L/940.pdf&Open=True} \href{https://docs.wto.org/dol2fe/Pages/SS/directdoc.aspx?filename=q\%3A/WT/L/940.pdf&Open=True}{ctdoc.aspx?filename=q:/WT/L/940.pdf\&O} \href{https://docs.wto.org/dol2fe/Pages/SS/directdoc.aspx?filename=q\%3A/WT/L/940.pdf&Open=True}{pen=True}

Moreover, the member exporting such a restriction must provide a notice in writing, as far in advance as possible, to the Committee on Agriculture, and consult, upon request, with any other member that is a prospective importer.\footnote{WTO. (1994d). Disciplines on Export Prohibitions and Restrictions (Article 12). Agreement on Agriculture. World Trade Organization. \url{https://www.wto.org/english/docs_e/legal_e/ag_e.htm} art12} This provision balances a country’s contemporary internal security concerns (domestic food supply to be specific) with its commitments to the free market, extensively alleviating panic-driven trade wars while keeping global agricultural commerce stable and functional.

\subsection{United Nations Security Council}

\emph{\textbf{Resolution 1540 (2004):}} United Nations Security Council (UNSC) Resolution 1540 is a binding instrument adopted under the Chapter VII of the UN Charter that obligates all states to block non-state actors from gaining, producing, or using any sorts of weapons of mass destruction (WMD), and their methods of delivery, by means of effective national legislation and controls. This resolution indirectly affects global trade by mandating that all UN member states devise and enforce rigorous domestic restraint over materials, equipments, that could be used in the production of nuclear, chemical, or biological weapons.\footnote{UNSC. (2004, April 28). Resolution 1540 (2004) [S/RES/1540 (2004)]. United Nations Digital Library. \url{https://digitallibrary.un.org/record/520326}}

The provisions of this exemplary instrument require countries to establish extensive export, transit, and transshipment regulations along with durable border security measures. While this fundamentally causes friction into the free movement of goods, the resolution serves as a globally accepted standard that rationalizes the imposition of trade compliance barriers. Resolution 1540 thus makes supply chain security an universal benchmark rather than an erratic trade barrier by legally obligating countries to strictly oversee and restrict dual-use supply chains.\footnote{UNSC. (2004, April 28). Resolution 1540 (2004) [S/RES/1540 (2004)]. United Nations Digital Library. \url{https://digitallibrary.un.org/record/520326}}

\subsection{United Nations Security Council}

\emph{\textbf{Resolution 2V1Y (2018):}} The catastrophic impact of armed conflict on global food supply chains is addressed by the UNSC Resolution 2417, which was adopted unanimously in 2018. Although the UN Security Council often authorizes economic sanctions, this particular resolution functions as a protective exception for agricultural commerce. Explicitly condemning starvation of civilians as a warfare tactic, it proscribes unlawful denial of humanitarian access to civilian populations, which essentially protects the logistics and infrastructure necessary to move food and agricultural goods.\footnote{UNSC. (2018, May 24). Resolution 2416 (2018) [S/RES/2417 (2018)]. UNDOCS. \url{https://docs.un.org/en/S/res/2417(2018)}}

The resolution’s provisions legally distinguish the trade of vital foodstuffs from legitimate military or economic targets. As per its mandate, even amid active hostilities and geopolitical blockades, states and combatants must grant safe, unhindered passage of food commodities and agricultural relief.\footnote{UNSC. (2018, May 24). Resolution 2416 (2018) [S/RES/2417 (2018)]. UNDOCS. \url{https://docs.un.org/en/S/res/2417(2018)}} This builds a legal shield that averts the complete shutdown of agricultural trade in conflict zones, ensuring the survival of the global food market while stabilizing military objectives.

\subsection{United Nations General Assembly}

\emph{\textbf{Resolution Y6/161 (Human Riehts and Unilateral Cogrcivg Mgasurgs, 2021):}} The UN General Assembly (UNGA) Resolution 76/161, titled \emph{Human Rights and Unilateral Coercive Measures}, categorically condemns the usage of unilateral coercive measures as means of political and economic pressure, delivering a crucial counterbalance to the weaponization of trade in an era of rising geopolitical tension. This resolution strongly urges all states to refrain from adopting, maintaining, or implementing any unilateral measures not in accordance with international law, international humanitarian law, and the UN Charter, specially highlighting those of a coercive nature with extraterritorial effects that hinder trade relations among states.\footnote{UNGA. (2021, December 16). Human rights and unilateral coercive measures [A/RES/76/161]. United Nations Digital Library. \url{https://digitallibrary.un.org/record/395494} 1} It anchors its legal stance in Article 32 of the Charter of Economic Rights and Duties of States, and formally rejects the use of economic, political, or any other type of measures to compel another state into the “subordination of the exercise of its sovereign rights”.\footnote{UNGA. (1974, December 12). Article 32. Charter of Economic Rights and Duties of States [Resolution 3281 (XXIX)]. UNDOCS. \url{https://docs.un.org/en/a/res/3281(XXIX)}} The resolution seeks to safeguard the global free trade system from being disintegrated by individual states that capitalize on broad, unverified national security claims to justify arbitrary economic warfare against geopolitical rivals.

Regarding the direct intersection of national security and free trade, A/RES/76/161 stresses that critical international trade, particularly involving food, agricultural products, medicine and medical equipment, and basic civilian commodities, must never be utilized as tools for exerting political and/or economic pressure. Furthermore, it explicitly condemns the inclusion of member states in unilateral lists under false pretexts of national security or terrorism, acknowledging that such listings haphazardly distort global trade and disrupt supply chains.\footnote{UNGA. (2021, December 16). Human rights and unilateral coercive measures [A/RES/76/161]. United Nations Digital Library. \url{https://digitallibrary.un.org/record/395494} 1} While it is non-binding in nature, this resolution serves as a normative mechanism to balance national security claims against open trade, by establishing the notion that even when nations adopt defensive security stances or engage in geopolitical disputes, open trade channels for life-saving goods are preserved, guaranteeing that national security strategies does not transform into indiscriminate obstructions to civilian commerce.

\emph{\textbf{SAFE Framework of Standards - World Customs Organization (WCO):}} Adopted in 2005 and updated regularly, the WCO SAFE Framework of Standards (SAFE FoS) was developed precisely in response to intensified global security concerns to safeguard the international supply chain without choking international commerce. Although it is not binding in nature, it provides an instrumental set of guidelines that customs administrations all across the world have integrated into domestic law.

At the core of the mechanism of the SAFE FoS is the Authorized Economic Operator (AEO) concept. It incentivizes the private sector to adopt meticulous internal security practices.\footnote{WCO. (2025, June). SAFE Framework of Standards. World Customs Organization.} Display of highly secure supply chains by the practicing companies can result in expedited customs processing and flexible border inspections. Such a mechanism functionally balances security and free trade through rewarding compliance with trade facilitation, and attempts to enable enhanced border security to target high-risk shipments while streamlining free flow of legitimate commerce.

\emph{\textbf{OECy Guidglings for Recipient Country Investment Policies Relating to National Security (2009):}} Since geopolitical tensions growingly manifest as foreign direct investment (FDI) in critical infrastructure and technology, the Organisation for Economic Co-operation and Development (OECD) addressed the necessity to balance open investment environments with national security screening. The Recommendation of the Council on Guidelines for Recipient Country Investment Policies relating to National Security, adopted by the OECD Council in 2009, set up a multilateral framework for how governments should review foreign investments to prevent malevolent geopolitical actors from obtaining vital domestic assets.\footnote{OECD. (2009, May 25). Recommendation of the Council on Guidelines for Recipient Country Investment Policies relating to National Security. OECD/LEGAL/0372. \url{https://legalinstruments.oecd.org/public/} doc/227/227.en.pdf}

The fundamental concepts of the Guidelines   are  centered  on non-discrimination, transparency, predictability, proportionality and accountability. It advises that FDI inspection frameworks should be closely tailored to genuine national security threats instead of acting as a covert tool

\url{https://www.wcoomd.org/-/media/wco/public/global/pdf/topics/facilitation/instruments-and-tools/tools/safe-package/safe-framework-2025_en.pdf?la=en} \href{https://www.wcoomd.org/-/media/wco/public/global/pdf/topics/facilitation/instruments-and-tools/tools/safe-package/safe-framework-2025_en.pdf?la=en}{blic/global/pdf/topics/facilitation/instrum} \href{https://www.wcoomd.org/-/media/wco/public/global/pdf/topics/facilitation/instruments-and-tools/tools/safe-package/safe-framework-2025_en.pdf?la=en}{ents-and-tools/tools/safe-package/safe-fr} \href{https://www.wcoomd.org/-/media/wco/public/global/pdf/topics/facilitation/instruments-and-tools/tools/safe-package/safe-framework-2025_en.pdf?la=en}{amework-2025\_en.pdf?la=en} for economic protectionism.\footnote{OECD. (2009, May 25). Recommendation of the Council on Guidelines for Recipient Country Investment Policies relating to National Security. OECD/LEGAL/0372. \url{https://legalinstruments.oecd.org/public/} doc/227/227.en.pdf} Thus, through standardizing the criteria for investment screening, the OECD framework assists in maintaining the free flow of global capital by avoiding national security debates from degrading into obscure impediments in investment.

\emph{\textbf{Thg Wassgnaar Arraneement (1996):}} The Wassenaar Arrangement on Export Controls for Conventional Arms and Dual-Use Goods and Technologies is the first global multilateral regime on export controls for conventional weapons and sensitive dual-use goods and technologies. Geopolitical tension linked to transfer of technology is addressed directly by this arrangement as its provisions aim to prevent disruptive accumulations of conventional arms and dual-use technologies. Nations participating in this arrangement consent to uphold national exports controls on mutually agreed lists of vastly sensitive technologies, which include but are not limited to advanced electronics, aerospace components, and cybersecurity tools.\footnote{Wassenaar Arrangement. (1996). The Wassenaar Arrangement on Export Controls for Conventional Arms and Dual-Use Goods and Technologies. \url{https://www.wassenaar.org/the-wassenaa} r-arrangement/}

The mechanism functions through sharing of information and voluntary compliance, where participating states integrate the Wassenaar control lists into their domestic export licensing systems.\footnote{Wassenaar Arrangement. (1996). The Wassenaar Arrangement on Export Controls for Conventional Arms and}  This  framework  adopts a

targeted approach rather than a blanket-ban on high-tech trade, letting the vast majority of consumer technology to trade freely while calling for rigorous, coordinated licensing requirements on strategic technologies.

\emph{\textbf{Article Y (Transit) of the Energy Charter Treaty (1994):}} One of the most robust legal mechanisms to prevent the weaponization of energy commerce during geopolitical disputes is provided by Article 7 of the Energy Charter Treaty. By acknowledging that energy safety is a crucial national interest, it explicitly prohibits contracting parties from disrupting (or allowing the disruption of) and reducing the existing flow of energy materials and goods (e.g. oil and gas via pipelines) at the time of an ongoing political or trade dispute.\footnote{Energy Charter Treaty. (1994). Article 7: Transit. energycharter.org. \url{https://www.energychartertreaty.org/} provisions/part-ii-commerce/article-7-transit/}

This provision serves as a buffer: before a country can pause energy transit due to a geopolitical grievance, it must resort to specific dispute resolution procedures highlighted in the treaty. Consequently, Article 7 helps to strike a balance between a nation’s right to resolve transit disputes and the global necessity of uninterrupted energy trade.\footnote{Energy Charter Treaty. (1994). Article 7: Transit. energycharter.org. \url{https://www.energychartertreaty.org/} provisions/part-ii-commerce/article-7-transit/} Ultimately, it aims to ensure that abrupt geopolitical escalations do not instantly cut-off the energy supplies of third-party countries down the supply chain.

Dual-Use Goods and Technologies. \url{https://www.wassenaar.org/the-wassenaar-arrangement/} \href{https://www.wassenaar.org/the-wassenaar-arrangement/}{r-arrangement/}

\subsection{Article 23 of the Geneva Convention}

\emph{\textbf{(IV) relative to the Protection of Civilian Persons in Time of War (1949):}} While mainly serving as the foundation of International Humanitarian Law, Article 23 of the Fourth Geneva Convention establishes an important basis for international commerce: free passage of shipments of medical supplies, food, and clothing, even during times of adverse geopolitical crisis.\footnote{International Committee of the Red Cross (ICRC). (1949). Consignment of medical supplies, food and clothing (Article 23). Convention (IV) relative to the Protection of Civilian Persons in Time of War. IHL Databases. \url{https://ihl-databases.icrc.org/en/ihl-treati} es/gciv-1949/article-23} This legally obligates all High Contracting Parties to allow the free passage of all consignments of medical and hospital stores, essential foodstuffs, clothing and tonics intended for children under fifteen, expectant mothers and maternity cases.

Article 23 requires states enforcing a military or economic blockade to inspect, approve, and safely route these specific commercial products through their lines.\footnote{International Committee of the Red Cross (ICRC). (1949). Consignment of medical supplies, food and clothing (Article 23). Convention (IV) relative to the Protection of Civilian Persons in Time of War. IHL Databases. \url{https://ihl-databases.icrc.org/en/ihl-treati} es/gciv-1949/article-23} Its mechanism compels belligerents to balance their maximum national security posture with the preservation of vital life, by depriving states from the legal right to impose absolute, impervious trade embargoes. Therefore, its ramification, albeit indirectly, is linked to saving commerce, and guaranteeing that trade in essential human goods cannot be completely blocked.
